\documentclass[a4paper,12pt,oneside]{maki-notes}

\renewcommand{\LectureNotesTitle}{Research Workflow Test}
\renewcommand{\AuthorInfo}{Test Runner}

\begin{document}
\frontmatter
\tableofcontents

\mainmatter

\SetResearchInfo{
  project={紧算子与谱结构},
  area={泛函分析},
  author={Maki},
  collaborators={A. B.},
  status={working draft},
  created={2026-04-13},
  updated={2026-04-13},
  tags={compact operator, spectrum}
}

\chapter{紧算子谱笔记}

\SetEntryInfo{
  topic={Riesz--Schauder 结构},
  source={Reed--Simon I},
  sourcekey={reed-simon-1},
  updated={2026-04-13},
  tags={compact operator, Fredholm},
  objective={整理非零谱离散性的证明骨架}
}

\MakeEntryGuide{
  summary={当前条目整理紧算子非零谱离散性的证明结构。},
  dependencies={Fredholm alternative; compact perturbation},
  targets={补全非零谱特征值有限重数部分},
  blockers={resolvent 紧性的引理链还不够紧凑},
  nextsteps={核对 Kato 的表述; 补 lemma 的引用位置}
}

\section{基础对象}

\DeclareSymbol{$\sigma(T)$}{算子 $T$ 的谱}
\DeclareSymbol{$\rho(T)$}{算子 $T$ 的 resolvent set}

\begin{notation}[谱与 resolvent]
记 $\sigma(T)$ 为算子 $T$ 的谱，$\rho(T)$ 为其 resolvent set。
\end{notation}

\begin{claim}[非零谱候选]
\label{clm:nonzero-spectrum}
若 $\lambda \neq 0$ 且 $\lambda \in \sigma(T)$，则它应表现出离散谱点的行为。
\end{claim}

\begin{remark}
这一节只整理论证骨架，不追求最强一般性。
\end{remark}

\begin{proofsketch}
先把 $(\lambda I - T)^{-1}$ 的分析拆成有限秩逼近与 Fredholm 备选两步。
\end{proofsketch}

\section{开放问题}

\begin{question}[更短证明链]
能否不调用最完整的 Fredholm alternative，而是直接证明非零谱只能由有限重特征值组成？
\end{question}

\begin{conjecture}[谱点有限重]
对 Banach 空间上的紧算子，所有非零谱点都对应有限代数重数。
\end{conjecture}

\begin{idea}[从 resolvent 展开入手]
先把 resolvent identity 写成可迭代的局部形式，再决定是否需要更强的 Fredholm machinery。
\end{idea}

\begin{obstacle}
目前缺少一个足够短的引理链把有限秩逼近与局部可逆性拼起来。
\end{obstacle}

\begin{nextstep}
把 Claim~\ref{clm:nonzero-spectrum} 改写为可独立引用的 lemma。
\end{nextstep}

\PrintSymbolIndex
\PrintOpenProblemIndex

\end{document}
